% !TEX root = ../main.tex
\corrigesChapitre{8}{Piles et oxydoréduction}
\corrigesSouples

\begin{blocLivre}{cor-8-1}
\begin{corrige}{1}{Écrire des demi-équations électroniques}
\demarche{Équilibrer les éléments, puis l'oxygène avec \ce{H2O}, l'hydrogène avec \ce{H+} et les charges avec des électrons.}
\begin{questions}
\item {\renewcommand{\arraystretch}{1.25}\begin{tabular}[t]{@{}l@{\hspace{3mm}}l@{\hspace{10mm}}l@{\hspace{3mm}}l@{}}
A. & \ce{Ag+ + e- <=> Ag} & D. & \ce{O2 + 4 H+ + 4 e- <=> 2 H2O}\\
B. & \ce{Fe^3+ + e- <=> Fe^2+} & E. & \ce{2 H+ + 2 e- <=> H2}\\
C. & \ce{I2 + 2 e- <=> 2 I-} & F. & \ce{MnO4- + 8 H+ + 5 e- <=> Mn^2+ + 4 H2O}
\end{tabular}}
\item Le cuivre cède des électrons : il est \res{\text{oxydé}}, c'est le \res{\text{réducteur}}. Les ions \ce{Ag+} captent des
électrons : ils sont \res{\text{réduits}}, ce sont les \res{\text{oxydants}}.
\item On combine \ce{Fe^2+ -> Fe^3+ + e-}, multipliée par 4, et \ce{O2 + 4 H+ + 4 e- -> 2 H2O} :
\[\res{\ce{4 Fe^2+ + O2 + 4 H+ -> 4 Fe^3+ + 2 H2O}}.\]
\item On combine \ce{Fe^2+ -> Fe^3+ + e-}, multipliée par 5, et \ce{MnO4- + 8 H+ + 5 e- -> Mn^2+ + 4 H2O} :
\[\res{\ce{MnO4- + 5 Fe^2+ + 8 H+ -> Mn^2+ + 5 Fe^3+ + 4 H2O}}.\]
\end{questions}
\end{corrige}
\end{blocLivre}
\begin{blocLivre}{cor-8-2}
\begin{corrige}{2}{L'arbre de Diane}
\demarche{Le cuivre cède des électrons aux ions argent directement, à la surface du fil.}
\begin{questions}
\item Les cristaux brillants ramifiés montrent la formation d'argent métallique ; la coloration bleue de la solution
révèle l'apparition d'ions \ce{Cu^2+}.
\item \ce{Cu -> Cu^2+ + 2 e-} et \ce{Ag+ + e- -> Ag} (multipliée par 2). Les ions \res{\ce{Ag+}\text{ sont l'oxydant}},
le \res{\text{cuivre est le réducteur}}.
\item $Q_\mathrm{r}=\dfrac{[\ce{Cu^2+}]}{[\ce{Ag+}]^2}$ ; au départ, il n'y a pas d'ions \ce{Cu^2+} :
$\res{Q_{\mathrm{r},\mathrm{i}}=0<K}$, le système évolue spontanément dans le sens direct.
\item $n(\ce{Ag+})=0{,}0500\times0{,}100=\qty{5,00e-3}{mol}$, d'où $m(\ce{Ag})=\num{5.00e-3}\times107{,}9\approx\res{\qty{0,540}{g}}$ ;
il se consomme $\qty{2,50e-3}{mol}$ de cuivre, soit $\res{\qty{0,159}{g}}$.
\item \res{\text{Non}} : les électrons passent directement du cuivre aux ions argent, sans circuler dans un circuit
extérieur. L'énergie libérée est dissipée sous forme de chaleur.
\end{questions}
\end{corrige}
\end{blocLivre}
\begin{blocLivre}{cor-8-3}
\begin{corrige}{3}{Déterminer la polarité d'une pile}
\demarche{Une tension positive indique que la borne V est reliée au pôle positif ; la réduction a lieu au pôle positif.}
\begin{questions}
\item La tension affichée est positive : la borne V, reliée à l'argent, est au pôle positif. L'électrode
\res{\text{d'argent est le pôle }+}, l'électrode \res{\text{de cuivre le pôle }-}.
\item À la cathode (argent, pôle $+$) : \ce{Ag+ + e- -> Ag}, réduction ; à l'anode (cuivre, pôle $-$) :
\ce{Cu -> Cu^2+ + 2 e-}, oxydation.
\item $\res{\ce{Cu(s) + 2 Ag^+(aq) -> Cu^2+(aq) + 2 Ag(s)}}$. Initialement,
$Q_{\mathrm{r},\mathrm{i}}=\dfrac{0{,}10}{0{,}10^2}=10<K$ : le système évolue dans le sens direct, ce qui correspond bien à la
polarité observée.
\item Les électrons vont du cuivre vers l'argent dans le circuit extérieur ; le courant circule en sens inverse. Dans le
pont salin, les ions \ce{K+} migrent vers la solution de nitrate d'argent (cathode), les ions \ce{NO3-} vers la
solution de nitrate de cuivre(II) (anode).
\end{questions}
\schema[corrige]{08-ex03-cor}
\end{corrige}
\end{blocLivre}
\begin{blocLivre}{cor-8-4}
\begin{corrige}{4}{Le rôle du pont salin}
\demarche{Dans un circuit fermé, le courant est porté par les électrons dans les métaux et par les ions dans les solutions.}
\begin{questions}
\item Dans les fils, le courant est dû au déplacement des \res{\text{électrons}} ; dans les solutions et le pont salin, au
déplacement des \res{\text{ions}}.
\item Sans pont salin, le circuit est \res{\text{ouvert}} : aucun courant ne peut circuler, les réactions aux électrodes
s'arrêtent.
\item Dans le bécher de zinc, des ions \ce{Zn^2+} apparaissent : des anions \ce{NO3-} y migrent pour compenser leur charge.
Dans le bécher de cuivre, des ions \ce{Cu^2+} disparaissent : des cations \ce{K+} y migrent. Ainsi, \res{\ce{K+}\text{ vers
la cathode}, \ce{NO3-}\text{ vers l'anode}}, et chaque solution reste électriquement neutre.
\item Les ions \ce{Cu^2+} captent directement les électrons du zinc, à sa surface : il se forme du cuivre, mais aucun
électron ne circule dans le circuit extérieur. Séparer les réactifs oblige les électrons à passer par le circuit
extérieur : c'est ce qui \res{\text{produit un courant utilisable}}.
\end{questions}
\end{corrige}
\end{blocLivre}
\begin{blocLivre}{cor-8-5}
\begin{corrige}{5}{Oxydants et réducteurs du quotidien}
\demarche{Un oxydant capte des électrons, un réducteur en cède ; repérer le rôle de chaque espèce dans la situation.}
\begin{questions}
\item a. \ce{ClO-} : \res{\text{oxydant}} ; b. \ce{C6H8O6} : \res{\text{réducteur}} ; c. \ce{O2} : \res{\text{oxydant}} (le fer, réducteur,
est oxydé) ; d. \ce{H2} : \res{\text{réducteur}}.
\item \ce{ClO- + 2 H+ + 2 e- <=> Cl- + H2O} ; \ce{C6H6O6 + 2 H+ + 2 e- <=> C6H8O6} ; \ce{O2 + 4 H+ + 4 e- <=> 2 H2O}.
\item L'acide ascorbique, réducteur, réagit avec le dioxygène de l'air à la place des molécules du fruit : il est
\res{\text{oxydé en premier}} et protège ainsi la pomme du brunissement.
\item \ce{ClO-} capte des électrons : c'est \res{\text{l'oxydant}} ; \ce{Cl-} en cède : c'est \res{\text{le réducteur}}. En
milieu acide, il se dégage du \res{\text{dichlore, gaz très toxique}} : il ne faut jamais acidifier l'eau de Javel.
\end{questions}
\end{corrige}
\end{blocLivre}
\begin{blocLivre}{cor-8-6}
\begin{corrige}{6}{La pile Daniell}
\demarche{La polarité découle du signe de la tension ; la capacité se calcule avec le réactif limitant, deux électrons étant échangés par ion \ce{Cu^2+}.}
\begin{questions}
\item La tension est positive et la borne COM est reliée au zinc : \res{\text{le zinc est le pôle }-}, le cuivre le pôle $+$.
Au pôle $-$ : \ce{Zn -> Zn^2+ + 2 e-} ; au pôle $+$ : \ce{Cu^2+ + 2 e- -> Cu}.
\item $Q_{\mathrm{r},\mathrm{i}}=\dfrac{[\ce{Zn^2+}]_\mathrm{i}}{[\ce{Cu^2+}]_\mathrm{i}}=\res{1{,}0}<K$ : sens direct, le zinc est oxydé et les ions
\ce{Cu^2+} sont réduits, ce qui correspond à la polarité observée.
\item $n(\ce{Zn})=\frac{5{,}0}{65{,}4}=\qty{7,6e-2}{mol}$ et $n(\ce{Cu^2+})=\qty{1,0e-2}{mol}$ : les \res{\text{ions }\ce{Cu^2+}\text{ sont
limitants}}. $Q=2\times\num{1.0e-2}\times\num{96500}=\res{\qty{1,93e3}{C}}$, soit $\frac{1930}{3{,}6}\approx\res{\qty{536}{mA.h}}$.
\item $\Delta t=\dfrac{Q}{I}=\dfrac{1930}{0{,}050}=\qty{3,86e4}{s}\approx\res{\qty{10,7}{h}}$.
\item Il se dépose $\qty{1,0e-2}{mol}$ de cuivre : $\Delta m=\num{1.0e-2}\times63{,}5\approx\res{\qty{0,64}{g}}$.
\end{questions}
\end{corrige}
\end{blocLivre}
\begin{blocLivre}{cor-8-7}
\begin{corrige}{7}{La pile d'une montre}
\demarche{L'oxyde d'argent est réduit : son électrode est la cathode. Comparer les quantités des deux réactifs, consommés mole à mole.}
\begin{questions}
\item L'oxyde d'argent capte des électrons, il est réduit : \res{\text{l'électrode à l'oxyde d'argent est le pôle }+}.
\item En additionnant les demi-équations : $\res{\ce{Zn + Ag2O -> ZnO + 2 Ag}}$. Les ions hydroxyde, consommés à l'anode
et régénérés à la cathode, \res{\text{n'apparaissent pas dans le bilan}}.
\item $n(\ce{Zn})=\frac{0{,}034}{65{,}4}=\qty{5,2e-4}{mol}$ et $n(\ce{Ag2O})=\frac{0{,}15}{231{,}7}=\qty{6,5e-4}{mol}$ : le
\res{\text{zinc est limitant}}. $Q=2\times\num{5.2e-4}\times\num{96500}\approx\qty{100}{C}$, soit $\res{\qty{28}{mA.h}}$.
\item $\Delta t=\dfrac{100}{\num{1.5e-6}}\approx\qty{6,7e7}{s}$, soit $\res{\text{environ 2 ans}}$.
\item Elles contiennent de l'argent, métal précieux à recycler, du zinc et une solution d'hydroxyde de potassium
corrosive : jetées avec les ordures ménagères, elles \res{\text{pollueraient l'environnement}}.
\end{questions}
\end{corrige}
\end{blocLivre}
\begin{blocLivre}{cor-8-8}
\begin{corrige}{8}{Suivre l'usure d'une pile}
\demarche{Pendant que la pile débite, deux électrons circulent pour chaque ion \ce{Cu^2+} réduit : $n(\mathrm{e}^-)=2\,\Delta n(\ce{Cu^2+})$.}
\begin{questions}
\item $\res{\ce{Cu^2+(aq) + Zn(s) -> Cu(s) + Zn^2+(aq)}}$.
\item Les points sont alignés : $[\ce{Cu^2+}]$ \res{\text{diminue linéairement}}. Le coefficient directeur de la droite
moyenne vaut $a\approx\res{\qty{-0,0187}{mol.L^{-1}.h^{-1}}}$.
\item $n(\mathrm{e}^-)=2V\,\big([\ce{Cu^2+}]_0-[\ce{Cu^2+}]_t\big)=2V\,|a|\,t$ et $n(\mathrm{e}^-)=\dfrac{I\,t}{F}$, d'où
\[I=2FV|a|=2\times\num{96500}\times0{,}100\times\frac{0{,}0187}{3600}\approx\res{\qty{0,10}{A}}.\]
\item La pile est usée lorsque les ions \ce{Cu^2+} sont épuisés : $t=\dfrac{0{,}100}{0{,}0187}\approx\res{\qty{5,3}{h}}$. Sa capacité
vaut $Q=2F\,n(\ce{Cu^2+})_0=2\times\num{96500}\times\num{1.00e-2}\approx\res{\qty{1,9e3}{C}}$, soit \qty{0,54}{A.h}.
\item \res{\text{Non}} : pour que $Q_\mathrm{r}=\frac{[\ce{Zn^2+}]}{[\ce{Cu^2+}]}$ atteigne $K$, il faudrait
$[\ce{Cu^2+}]\approx\frac{0{,}200}{\num{1.9e37}}\approx\qty{1e-38}{mol.L^{-1}}$. La pile s'arrête parce que le réactif limitant,
\ce{Cu^2+}, a disparu : la transformation est totale.
\end{questions}
\schema[corrige]{08-ex08}
\end{corrige}
\end{blocLivre}
\begin{blocLivre}{cor-8-9}
\begin{corrige}{9}{La pile zinc-air des aides auditives}
\demarche{Convertir la capacité ($\qty{1}{mA.h}=\qty{3,6}{C}$), puis utiliser deux électrons par atome de zinc oxydé.}
\begin{questions}
\item Le dioxygène capte des électrons : c'est \res{\text{l'oxydant}} ; le zinc, qui en cède, est \res{\text{le réducteur}}.
Le zinc est oxydé : son électrode est l'anode, \res{\text{le pôle }-}.
\item En multipliant la première demi-équation par 2 : $\res{\ce{2 Zn + O2 -> 2 ZnO}}$.
\item $Q=620\times3{,}6=\res{\qty{2232}{C}}$ ; $n(\mathrm{e}^-)=\frac{2232}{96500}=\qty{2,31e-2}{mol}$, donc
$n(\ce{Zn})=\qty{1,16e-2}{mol}$ et $m(\ce{Zn})\approx\res{\qty{0,76}{g}}$.
\item La pile fournit $\frac{620}{1{,}6}\approx\qty{388}{h}$ de fonctionnement, soit $\frac{388}{16}\approx\res{24\text{ jours}}$.
\item L'oxydant, le dioxygène, n'est pas stocké dans la pile : il vient de l'air, ce qui \res{\text{allège la pile}}. Une
fois l'opercule retiré, le dioxygène entre et la pile commence à s'user, même si l'appareil ne fonctionne pas.
\end{questions}
\end{corrige}
\end{blocLivre}
\begin{blocLivre}{cor-8-10}
\begin{corrige}{10}{Le sodium, un métal très réducteur}
\demarche{Un métal du bloc s cède ses électrons de valence pour acquérir la configuration du gaz noble qui le précède.}
\begin{questions}
\item $1s^2\,2s^2\,2p^6\,3s^1$ : le sodium possède \res{\text{un électron de valence}}.
\item En perdant cet électron, il acquiert la configuration stable du néon : il cède très facilement un électron, c'est
un réducteur puissant. Il forme \res{\text{l'ion }\ce{Na+}}.
\item On combine \ce{Na -> Na+ + e-}, multipliée par 2, et \ce{2 H2O + 2 e- -> H2 + 2 HO-} :
\[\res{\ce{2 Na + 2 H2O -> 2 Na+ + 2 HO- + H2}}.\]
Les ions hydroxyde formés rendent la solution basique : la phénolphtaléine devient rose.
\item $n(\ce{Na})=\frac{0{,}23}{23{,}0}=\qty{1,0e-2}{mol}$, d'où $n(\ce{H2})=\qty{5,0e-3}{mol}$ et
$V=\num{5.0e-3}\times24\approx\res{\qty{0,12}{L}}$.
\item L'huile isole le sodium de l'eau et du dioxygène de l'air, avec lesquels il réagirait. \ce{Mg^2+ + 2 e- <=> Mg} : comme
le sodium, le magnésium cède facilement ses électrons de valence (deux) ; c'est \res{\text{un réducteur}}.
\end{questions}
\end{corrige}
\end{blocLivre}
\begin{blocLivre}{cor-8-11}
\begin{corrige}{11}{Simuler la décharge d'une pile}
\demarche{À intensité constante, l'avancement augmente de $I\,\mathrm{d}t/(2F)$ à chaque pas de temps.}
\begin{questions}
\item \ce{Cu^2+ + Zn -> Cu + Zn^2+}. \code{n\_cu} et \code{n\_zn} sont les quantités d'ions \ce{Cu^2+} et \ce{Zn^2+}
présentes à la date \code{t}.
\item Pendant \code{dt}, la charge $I\,\mathrm{d}t$ circule, soit $\frac{I\,\mathrm{d}t}{F}$ mol d'électrons ; deux électrons étant
échangés pour un avancement de 1~mol, l'avancement augmente de $\res{\frac{I\,\mathrm{d}t}{2F}}$.
\item La boucle s'arrête lorsque les ions \ce{Cu^2+} sont épuisés. Le programme affiche la durée de fonctionnement
(\res{\qty{19320}{s}}, soit 5,37~h) et la quantité finale d'ions \ce{Zn^2+} (\qty{0,010}{mol}, le double de la
quantité initiale).
\item $Q=2F\,n(\ce{Cu^2+})_0=2\times\num{96500}\times\num{5.0e-3}=\qty{965}{C}$, d'où $\Delta t=\frac{965}{0{,}050}=\res{\qty{19300}{s}}$.
Le programme avance par pas de 60~s : il s'arrête au premier multiple de 60~s qui dépasse cette durée.
\item L'intensité ne reste pas constante en fin de décharge : quand les ions \ce{Cu^2+} se raréfient, la résistance
interne de la pile augmente et le courant diminue.
\end{questions}
\end{corrige}
\end{blocLivre}
\begin{blocLivre}{cor-8-12}
\begin{corrige}{12}{La pile d'un stimulateur cardiaque}
\demarche{Capacité en coulombs, quantité d'électrons, puis masses : un électron par atome de lithium, deux par molécule de diiode.}
\begin{questions}
\item $1s^2\,2s^1$ : le lithium cède facilement son unique électron de valence pour acquérir la configuration de
l'hélium, c'est un réducteur puissant. \ce{Li+ + e- <=> Li} et \ce{I2 + 2 e- <=> 2 I-}.
\item Le lithium est oxydé : c'est \res{\text{l'anode, pôle }-} ; le diiode est réduit à la \res{\text{cathode, pôle }+}.
Équation de fonctionnement : $\res{\ce{2 Li + I2 -> 2 LiI}}$.
\item Le lithium, métal du bloc s, \res{\text{réagirait avec l'eau}} en libérant du dihydrogène : la pile serait détruite.
\item $Q=\qty{1,0}{A.h}=\qty{3600}{C}$ et $\Delta t=\frac{3600}{\num{1.0e-5}}=\qty{3,6e8}{s}$, soit $\res{\text{environ 11 ans}}$.
$n(\mathrm{e}^-)=\frac{3600}{96500}=\qty{3,73e-2}{mol}$ : $m(\ce{Li})=\num{3.73e-2}\times6{,}94\approx\res{\qty{0,26}{g}}$ et
$m(\ce{I2})=\frac{\num{3.73e-2}}{2}\times253{,}8\approx\res{\qty{4,7}{g}}$.
\item Le zinc échange deux électrons par atome : $m(\ce{Zn})=\frac{\num{3.73e-2}}{2}\times65{,}4\approx\res{\qty{1,2}{g}}$, soit près de
cinq fois plus que de lithium. Léger et très réducteur, le lithium donne une pile compacte et de tension élevée.
\end{questions}
\end{corrige}
\end{blocLivre}
\begin{blocLivre}{cor-8-13}
\begin{corrige}{13}{Une voiture à hydrogène}
\demarche{Deux électrons par molécule de dihydrogène ; l'énergie électrique vaut $E=Q\,U$ pour l'ensemble des cellules.}
\begin{questions}
\item Le dihydrogène est oxydé sur l'électrode A : c'est \res{\text{l'anode}} ; le dioxygène est réduit sur l'électrode B,
\res{\text{la cathode}}. Les électrons vont de A vers B dans le circuit extérieur ; les ions \ce{H+}, produits en A et
consommés en B, traversent la membrane de A vers B.
\item $\res{\ce{2 H2 + O2 -> 2 H2O}}$. La membrane sépare les réactifs et laisse passer les ions : elle joue le rôle du
\res{\text{pont salin}}.
\item $n(\ce{H2})=\frac{5600}{2{,}0}=\qty{2,8e3}{mol}$, $n(\mathrm{e}^-)=\qty{5,6e3}{mol}$ et
$Q=\num{5.6e3}\times\num{96500}\approx\res{\qty{5,4e8}{C}}$.
\item $E=Q\,U=\num{5.4e8}\times0{,}70\approx\qty{3,8e8}{J}$, soit $\res{\qty{105}{kW.h}}$. Pour 650~km :
$\frac{105}{6{,}5}\approx\res{\qty{16}{kW.h}}$ pour 100~km, une consommation comparable à celle d'une voiture électrique à
batterie.
\item Il se forme autant de molécules d'eau que de molécules de dihydrogène consommées :
$m(\ce{H2O})=\num{2.8e3}\times18{,}0\approx\res{\qty{50}{kg}}$.
\end{questions}
\schema[corrige]{08-ex13}
\end{corrige}
\end{blocLivre}
\begin{blocLivre}{cor-8-14}
\begin{corrige}{14}{Une pile de concentration}
\demarche{Le même couple intervient aux deux électrodes : seules les concentrations diffèrent, et la pile fonctionne jusqu'à ce qu'elles s'égalisent.}
\begin{questions}
\item La tension est positive et la borne COM est reliée au compartiment 1 : \res{\text{la lame 1 est le pôle }-}, la lame 2
le pôle $+$. Compartiment 1 (anode) : \ce{Cu -> Cu^2+ + 2 e-} ; compartiment 2 (cathode) : \ce{Cu^2+ + 2 e- -> Cu}.
\item $\ce{Cu^2+}_{(2)}+\ce{Cu}_{(1)}\longrightarrow\ce{Cu}_{(2)}+\ce{Cu^2+}_{(1)}$ et
$Q_\mathrm{r}=\dfrac{[\ce{Cu^2+}]_1}{[\ce{Cu^2+}]_2}$. À l'équilibre, les deux solutions seraient identiques :
$[\ce{Cu^2+}]_1=[\ce{Cu^2+}]_2$, d'où \res{K=1}.
\item $Q_{\mathrm{r},\mathrm{i}}=\frac{\num{1.0e-2}}{0{,}10}=0{,}10<K$ : le système évolue dans le sens direct, le cuivre est oxydé
dans le compartiment 1, qui est bien le pôle $-$.
\item La pile est usée quand $Q_\mathrm{r}=K$ : avec l'avancement volumique $x$, $0{,}010+x=0{,}10-x$, soit
$x=\qty{0,045}{mol.L^{-1}}$ : $\res{[\ce{Cu^2+}]_1=[\ce{Cu^2+}]_2=\qty{0,055}{mol.L^{-1}}}$. Alors
$n(\mathrm{e}^-)=2\times0{,}045\times0{,}100=\qty{9,0e-3}{mol}$ et $Q\approx\res{\qty{8,7e2}{C}}$. La lame 1 perd
$\qty{4,5e-3}{mol}$ de cuivre, soit 0,29~g : ses 2,0~g suffisent.
\item La pile Daniell s'arrête lorsque son réactif limitant est épuisé, sa constante étant immense. Ici, $K=1$ : la pile
s'arrête \res{\text{à l'équilibre}}, alors que les deux réactifs sont encore présents.
\end{questions}
\end{corrige}
\end{blocLivre}
\begin{blocLivre}{cor-8-15}
\begin{corrige}{15}{Classer des métaux à l'aide de piles}
\demarche{Une tension positive signifie que la borne COM est au pôle négatif ; une tension négative, qu'elle est au pôle positif.}
\begin{questions}
\item Zinc -- cuivre : zinc $-$, cuivre $+$. Cuivre -- argent : COM sur l'argent et tension négative, donc argent $+$,
cuivre $-$. Fer -- cuivre : cuivre $+$, fer $-$.
\item Pôle $-$ : \ce{Zn -> Zn^2+ + 2 e-} ; pôle $+$ : \ce{Cu^2+ + 2 e- -> Cu} ; $\res{\ce{Zn + Cu^2+ -> Zn^2+ + Cu}}$.
\item Le cuivre, pôle $+$, est \qty{1,10}{V} au-dessus du zinc et \qty{0,78}{V} au-dessus du fer ; l'argent, pôle $+$, est
\qty{0,43}{V} au-dessus du cuivre : \res{\text{zinc à }\qty{-1,10}{V}}, \res{\text{fer à }\qty{-0,78}{V}},
\res{\text{argent à }\qty{+0,43}{V}} (voir l'échelle).
\item Le pôle négatif est toujours le plus bas : le pouvoir réducteur augmente vers le bas de l'échelle,
$\res{\text{argent}<\text{cuivre}<\text{fer}<\text{zinc}}$. Zinc -- argent : $0{,}43-(-1{,}10)=\res{\qty{1,53}{V}}$, zinc $-$ ;
zinc -- fer : $-0{,}78-(-1{,}10)=\res{\qty{0,32}{V}}$, zinc $-$.
\item \res{\text{Oui}} : le fer, plus réducteur que le cuivre, réduit les ions \ce{Cu^2+} :
\ce{Fe(s) + Cu^2+(aq) -> Fe^2+(aq) + Cu(s)}.
\end{questions}
\schema[corrige]{08-ex15-cor}
\end{corrige}
\end{blocLivre}
